% Section 5: Empirical Validation
% IEEE Computer — "Silent Decisions Are Type Errors"
% Authors: Soares et al., 2026
%
% Benchmark data collected with:
%   cargo bench --bench verdict_construction
% on the hardware described in Section 5.1.
% Raw output preserved in bench_results.txt (repository root).

\section{Empirical Validation}
\label{sec:benchmarks}

The Constitutional Enclosure Theorem guarantees that silent decisions
are compile-time errors.  A natural question follows: does the
compile-time guarantee impose a \emph{runtime cost}?  This section
answers that question empirically.

% ----------------------------------------------------------------
\subsection{Experimental Setup}
\label{sec:setup}

\paragraph{Hardware and software.}
All measurements were collected on a single core of an
Intel Xeon processor at 2.10~GHz (4 logical cores available, benchmarks
single-threaded), running Linux 6.18.5 x86\_64.
The implementation was compiled with \texttt{rustc}~1.93.1
(release profile: \texttt{opt-level~=~3}, LTO disabled).
Criterion~0.5.1 was used as the benchmarking harness.

\paragraph{Methodology.}
We used Criterion's \texttt{iter\_batched} API with
\texttt{BatchSize::SmallInput}.  This separates the \emph{setup
phase} (allocating the context buffer, cloning string arguments) from
the \emph{measurement phase} (constructing the \texttt{Verdict}),
ensuring that heap allocation latency is not attributed to the
type-system mechanisms under study.

To prevent the LLVM optimiser from caching the BLAKE3 hash across
iterations (which would underestimate the hashing cost), a one-byte
seed is incremented on each setup call, producing a distinct input
for every iteration.  The explanation string, jurisdiction, and
policy version are cloned from pre-allocated originals; their cost
is excluded from the measurement window.

\texttt{black\_box} wrappers are applied to inputs passed to
\texttt{EvidenceToken::new}, preventing constant-folding of the hash
computation.

Criterion reports a 95\% confidence interval over 100 samples per
benchmark; we quote the point estimate (mean) with the lower and
upper bounds of the interval.

% ----------------------------------------------------------------
\subsection{Results}
\label{sec:results}

\paragraph{Memory footprint.}
\begin{center}
\texttt{std::mem::size\_of::\textless{}Verdict\textgreater{}() = 152 bytes} (x86-64, rustc~1.93.1, Linux)
\end{center}
This fixed footprint comprises the 32-byte \texttt{Blake3Hash}, the
\texttt{ComplianceToken} fields (two heap-allocated strings and a
\texttt{u32}), the \texttt{Decision} discriminant, the explanation
\texttt{String}, and the 32-byte HMAC integrity seal.  The serialised
on-wire footprint (JSON or CBOR) is approximately 300--500 bytes
depending on field lengths.

\paragraph{Verdict construction latency.}
Table~\ref{tab:new} reports the end-to-end latency of
\texttt{Verdict::new}, which includes BLAKE3 context hashing,
HMAC-SHA256 sealing, and memory allocation for the \texttt{Verdict}
struct.  Three context sizes are measured to characterise the
BLAKE3 throughput curve.

\begin{table}[t]
\caption{Latency of \texttt{Verdict::new} by context size (95\% CI)}
\label{tab:new}
\begin{tabular}{lrrr}
\hline
\textbf{Context size} & \textbf{Lower} & \textbf{Mean} & \textbf{Upper} \\
\hline
64~B   & 414.67~ns & 416.80~ns & 419.08~ns \\
512~B  & 859.22~ns & 869.26~ns & 880.74~ns \\
4~KiB  & 1.639~\textmu{}s & 1.671~\textmu{}s & 1.708~\textmu{}s \\
\hline
\end{tabular}
\end{table}

\paragraph{Integrity verification latency.}
Table~\ref{tab:verify} reports the latency of
\texttt{Verdict::verify\_integrity}, which recomputes the
HMAC-SHA256 seal over the stored evidence identifier, decision, and
explanation and compares it to the seal embedded in the
\texttt{Verdict}.  This operation does not re-hash the original
context.

\begin{table}[t]
\caption{Latency of \texttt{Verdict::verify\_integrity} (95\% CI)}
\label{tab:verify}
\begin{tabular}{lrrr}
\hline
\textbf{Operation} & \textbf{Lower} & \textbf{Mean} & \textbf{Upper} \\
\hline
\texttt{verify\_integrity} & 324.09~ns & 327.36~ns & 330.32~ns \\
\hline
\end{tabular}
\end{table}

% ----------------------------------------------------------------
\subsection{Discussion}
\label{sec:discussion}

\paragraph{Zero-cost abstractions.}
The linear type discipline---private fields, \texttt{pub(crate)}
visibility, and \texttt{\#[must\_use]}---operates entirely at compile
time.  It introduces no additional branches, allocations, or
synchronisation primitives.  The runtime cost of a
\texttt{Verdict} is therefore identical to the cost of the
cryptographic operations it encapsulates (BLAKE3 and HMAC-SHA256),
which would be required regardless of the governance mechanism chosen.

\paragraph{Isolating the BLAKE3 cost.}
The difference between \texttt{Verdict::new} (64~B context) and
\texttt{Verdict::verify\_integrity} isolates the cost attributable
to hashing the decision context:
\[
  \Delta_{64\mathrm{B}} = 416.80\;\text{ns} - 327.36\;\text{ns}
                        \approx 89\;\text{ns}.
\]
This $\approx 89$~ns represents one BLAKE3 invocation over 64 bytes,
consistent with BLAKE3's documented throughput of approximately
1~GB/s on this hardware class.  As context size grows, the
additional latency scales proportionally: 512~B adds approximately
$452\;\text{ns}$ above the 64~B baseline, and 4~KiB adds
approximately $1.25\;\mu\text{s}$.

\paragraph{Governance at microsecond scale.}
All measured operations complete in under 2~\textmu{}s.  In a
production AI inference pipeline, where model evaluation typically
requires tens to hundreds of milliseconds, the governance overhead
introduced by the Constitutional Enclosure is
\emph{below measurement noise}---less than 0.01\% of total request
latency for any non-trivial model.  The accountability guarantee
is therefore practically free.

\paragraph{Verification cost.}
The \texttt{verify\_integrity} operation ($\approx 327$~ns) enables
an auditor to confirm the integrity of any stored verdict at a rate
exceeding $3 \times 10^6$ verifications per second on a single core.
This makes retroactive audit of large decision logs computationally
tractable without specialised hardware.

% ----------------------------------------------------------------
\subsection{Worked Examples}
\label{sec:worked-examples}

Three worked examples demonstrate BTV in semi-realistic scenarios: credit
scoring under LGPD (30-day appeal window), hiring screening under EU AI Act
Art.~86 (high-risk classification), and emergency triage (where \texttt{Allow}
decisions also require evidence).  Each example constructs verdicts from
domain-specific payloads and verifies integrity.  Source code is included in
the companion repository (\texttt{examples/}).
